\chapter{Las fuentes}
\label{cap:fuentes}

Este libro no tiene bibliografía en el sentido habitual, y conviene decir por qué antes de
enumerar nada. \textbf{Casi todo lo que afirma se apoya en fuentes primarias}: boletines
oficiales, expedientes administrativos y judiciales, resoluciones, decretos, planos y series
estadísticas. Discute literatura sólo de manera puntual y declarada: lo que hizo, sobre todo, fue leer documentos.

Ese apoyo se distribuye de manera desigual. El inventario normativo está en el apéndice
\ref{cap:normativa} y lo que falta pedir, ordenado por organismo, en el apéndice
\ref{cap:pedidos}. Lo que sigue es distinto: es el mapa de \textbf{qué clase de fuente
sostiene qué parte del libro}, para que el lector pueda calibrar cuánto pesa cada
afirmación según de dónde viene.

\section{Fuente seriada principal}

\textbf{Boletín Oficial de la Provincia de Salta}, desde su número 1 del 15 de octubre de
1908 hasta las ediciones de 2026. Es la columna vertebral del relevamiento y la fuente de la
mayor parte de las fechas, montos, resoluciones y nombres que el libro consigna. Se accedió
por el repositorio digital de la provincia; la búsqueda se hizo sobre texto extraído de los
PDF, de modo que \textbf{la cobertura no es exhaustiva}: un instrumento cuya digitalización
sea defectuosa no aparece, y su ausencia en este libro no prueba su inexistencia.

\textbf{Con una excepción declarada, y conviene que el lector la tenga presente porque cambia
lo que se puede afirmar.} Treinta y seis tramos se relevaron de otro modo: edición por edición, sin depender del término de
búsqueda. Treinta y dos de ellos ---que forman dos ventanas, 1908--1921 y 1923--1946, con el primer semestre de 1922 leído entre las dos--- se verificaron además una por una; los cuatro restantes, \textbf{1947 a 1950}, tienen pre-informe automático y no lectura sobre la imagen, y por eso quedan fuera de las ventanas. El primero va del \textbf{número 1, del 15 de octubre de
1908, al número 196, de octubre de 1910} ---las \textbf{196 primeras ediciones} de la historia
del Boletín, de las que faltan en el escaneo cuatro páginas impresas de 1908, la edición Nº 162 ---que el repositorio reemplaza por un aviso de que no la encontró--- y la última hoja de la Nº 157, las dos de 1910---. \textbf{La primera lectura de ese tramo no fue exhaustiva}: en septiembre de 2026 una segunda lectura de 1908, con OCR propio sobre las 103 hojas y transcripción sobre la imagen, encontró dos cosas que la primera había pasado por alto ---el comisario de policía y el clasificador de patentes del departamento, y los dueños del terreno del cateo de 1908--- y una que había leído de más: la ubicación del pleito de Wierna (capítulo \ref{cap:fincas}). En el mismo mes se releyó 1909 con ese procedimiento ---las 98 ediciones, 398 hojas--- y aparecieron veintidós actos con mención del departamento, varios de ellos ausentes de la primera lectura: el decreto que designa la comisión municipal, los veedores del Código Rural, los jueces de paz, la licitación del edificio para la policía y las tres renuncias de 1909. Uno de esos actos no lo encontró ninguna de las cuatro búsquedas automáticas, porque el corte de renglón partía la palabra con signos en medio, y apareció al buscar por nombres de personas. También se releyó 1910 ---98 ediciones, 391 hojas---, con veintiuna menciones del departamento, dos de ellas invisibles para las búsquedas automáticas porque el original parte la palabra con espacios o con guion blando. \textbf{La primera ventana quedó así releída entera}, y de esa relectura salen la comisión municipal de 1909 y 1910, los veedores del Código Rural, los doce partidos de junio de 1910, el contrato del edificio de la comisaría y los dos huecos de escaneo recién declarados. Lo que la primera lectura afirmó como ausente en 1908--1910 sigue en pie donde la relectura no dijo lo contrario; donde lo dijo, este libro lo corrigió. El segundo, del \textbf{número 962 al 2.464}, junio de 1923 a diciembre de 1945:
\textbf{1.503 números}, de los que faltan treinta y uno: \textbf{1.472 leídos}. Los faltantes
son el tramo del 977 al 990 (fines de septiembre a fines de diciembre de 1923); el del 1.037 al
1.042 (21 de noviembre a 26 de diciembre de 1924); el 1.095 (1 de enero de 1926); el 1.146 (24 de
diciembre de 1926); el 1.305 (10 de enero de 1930), que no está en el repositorio, y el 1.312,
1.318, 1.320, 1.323 y 1.327 (febrero a junio de 1930), que están pero en blanco; el 1.356 (2 de
enero de 1931); y el 1.512 y el 1.564, que \textbf{el índice del propio repositorio no registra}: la lista de 1933
termina en el Nº 1.511, del 22 de diciembre, y la de 1934 empieza en el 1.513, del 5 de enero; la de
1934 termina en el 1.563, del 21 de diciembre, y la de 1935 empieza en el 1.565, del 4 de enero. En
los dos casos falta el número del último viernes del año, y no se sabe si se publicó y no se
digitalizó o si la numeración saltó. \textbf{Los faltantes tienen, además, una estación}: siete de
los treinta y uno ---el 990 y el 1.042, que cierran sus tramos, y el 1.095, el 1.146, el 1.356, el
1.512 y el 1.564--- son números de la última semana de diciembre o la primera de enero. A eso se suman lagunas de página dentro de ediciones leídas: las páginas
impresas 49 a 51 del Nº 2.043 (1944) y hasta seis páginas de las ediciones de noviembre y
diciembre de 1945. El tercero, agregado al cierre, es \textbf{el año 1946 completo}: del número 2.465 (2 de enero) al 2.744 (31 de diciembre), \textbf{280 ediciones, todas presentes}, 3.912 hojas. Se leyó en septiembre de 2026 con otro procedimiento: OCR propio sobre todas las hojas, porque la capa de texto del Boletín venía con columnas cortadas, y lectura sobre la imagen de toda hoja con resultados, de las tablas y de las cifras citadas. Faltan en el escaneo 91 páginas, casi todas la última página par de su edición. El cuarto, intercalado entre el primero y el segundo, son \textbf{los años 1911 a 1914 enteros}. Los tres primeros: las ediciones 220 a 306, 347 hojas, en septiembre de 2026, de las que faltan en el escaneo tres hojas de la Nº 242, una de la Nº 243 y cuatro de la Nº 269; las ediciones 307 a 383, 307 hojas, de las que faltan la última hoja de la Nº 374 y la de la Nº 375; y las ediciones 384 a 458, 288 hojas de las 296 del año ---faltan completas la Nº 438 y la Nº 444---, con tres escaneos repetidos que el folio impreso permitió descontar. Los tres se leyeron con una diferencia de método que conviene declarar: esos años se leyeron sobre una sola capa de texto, la del propio Boletín, sin OCR propio que la contraste, y se miró en la imagen cada mención y las columnas de control ---veintiséis imágenes en 1911 y setenta en 1912; de 1913 no quedó asentado el recuento---. \textbf{El año 1914} se leyó en septiembre de 2026 con el procedimiento de 1946 y no con el de 1911--1913: las ediciones \textbf{459 a 536}, \textbf{279 hojas}, con \textbf{OCR propio sobre todas} ---sólo 23 de las ediciones traen capa de texto previa--- y lectura sobre la imagen de toda hoja con resultado. De las 78 ediciones nominales del año faltan en el repositorio las \textbf{487 a 493}, siete números consecutivos de junio, y la \textbf{470}, que el repositorio reemplaza por un aviso de que no la encontró: ocho ausencias sobre setenta y ocho, de modo que quedan \textbf{70 ediciones leídas} ---la cuenta se rehízo sobre el rango de ediciones y las ausencias declaradas, y corrige el 69 que traía el informe de barrido---. Las tapas de enero y febrero imprimen el año equivocado ---la Nº 461 dice «Enero 10 de 1913»--- y el calendario lo resuelve: el Boletín salía miércoles y sábados, y esas fechas caen así en 1914 y no en 1913. En 1912, tres actos ---el nombramiento de los jueces de paz, la integración de la comisión municipal y una entrega de la sentencia de Mojotoro--- no los devolvía ninguna búsqueda porque el original parte la palabra con guion blando y sin guion visible, y un cuarto ---el deslinde de Santa Mónica y Quitilipe--- porque la capa de texto lee «Oaldfera». \textbf{Y el quinto, agregado al cierre, es el año 1915}: las ediciones \textbf{537 a 586}, \textbf{195 hojas}, leídas en septiembre de 2026 con el procedimiento de 1914 y de 1946 ---OCR propio con modelo español sobre las 195 hojas, porque sólo diecinueve traen capa de texto previa, y lectura sobre la imagen de las cuarenta y nueve tapas y de las hojas con resultado---. \textbf{Es el tramo con más huecos declarados de los cinco}: la edición \textbf{581} no está en el repositorio y no hay ediciones de enero ni de la primera semana de febrero ---la serie empieza el 10 de febrero---, de modo que son \textbf{cuarenta y nueve ediciones leídas}. El barrido dejó trece apariciones del departamento en doce hojas, \textbf{ningún acto de tierra}, \textbf{ningún balance municipal} en todo el año y \textbf{ningún plano del departamento} entre los veintisiete que el año cita sin acompañar. Tres encabezados y la segunda publicación del edicto sucesorio se leyeron sobre el texto reconocido y no sobre la imagen, y así quedan declarados. \textbf{Y el sexto es el año 1916}: las ediciones \textbf{587 a 635}, \textbf{196 hojas}, \textbf{sin un solo salto de numeración} y con las cuatro hojas completas en las cuarenta y nueve. Leído en septiembre de 2026 con el mismo procedimiento: OCR propio con modelo español sobre las 196 hojas ---la capa previa cubre 92 y se concentra a partir de la edición 615---, lectura sobre la imagen de las cuarenta y nueve tapas y de veintidós hojas, y treinta y una columnas recortadas y leídas enteras. \textbf{Es el tramo más limpio de los seis}: no falta ninguna edición ni ninguna hoja. Dejó doce actos del departamento en catorce hojas, entre ellos \textbf{tres actos de dominio} que el apéndice \ref{cap:dominio} incorpora, y \textbf{ningún cuadro, ningún resumen de Tesorería y ningún plano}. Ciento setenta y cuatro hojas no se miraron en la imagen: ninguna de ellas dio acierto del término en ninguna de las cuatro pasadas, y así se declara, y no como ausencia de columna perdida en el año. \textbf{Y el séptimo es el año 1917}, y es el más flojo de los once primeros: las ediciones \textbf{636 a 683}, de las que \textbf{falta entera la 640}, con \textbf{289 hojas de las 329} que la paginación impresa declara para el año. Se leyó en septiembre de 2026 \textbf{sin OCR propio} ---el entorno no tenía el modelo español--- de modo que la única lectura posible fuera de la capa de texto del PDF fue la imagen: 76 hojas miradas, 112 recortes, y las 213 restantes sólo por la capa, que en este año está a unos 150 ppi, la mitad que en 1912. \textbf{El año se barrió dos veces, de manera independiente y con dos juegos distintos de términos} ---el segundo con ochenta y cuatro contra setenta y seis, y con doscientos treinta y tres actores contra doscientos trece---, y los dos barridos llegaron a \textbf{los mismos nueve actos}, uno mirando setenta y seis hojas en la imagen y el otro treinta y una. Es la única comprobación cruzada de este tipo que tiene el libro, y vale sobre todo en un tramo sin OCR propio. Dejó nueve actos del departamento en once hojas, \textbf{tres de ellos deslindes} que el apéndice \ref{cap:dominio} incorpora, y \textbf{ningún plano}: los tres deslindes se describen por linderos y rumbos, sin plano ni croquis, al revés que en 1912. \textbf{En este tramo, más que en ningún otro, un cero puede ser de la búsqueda y no del archivo}: el propio barrido nombra siete personas del departamento cuya búsqueda por nombre no devuelve ni su propio acto, porque la capa imprime «eres.sart», «keyes» y «olia ile». \textbf{Con esa reserva, hay un cero que conviene registrar}: de los ochenta y cuatro términos buscados, \textbf{cincuenta y cuatro no dieron un solo acierto en todo 1917}, y entre ellos están San Alejo, Santa Rufina, Los Porongos, Campo Alegre, Los Peñones, Potrero de Castillo y El Durazno. Las fincas grandes del departamento no aparecen en el año: los tres únicos actos de tierra recaen sobre la Calderilla y sobre el Potrero Gallinato. \textbf{Y el octavo es el año 1918}, agregado al cierre: las ediciones \textbf{684 a 736}, \textbf{534 hojas} sobre los \textbf{587 folios} que la paginación anual corrida declara para el año, de modo que faltan 53 hojas; \textbf{falta entera la edición 703}, del 17 de mayo, que el repositorio reemplaza por un aviso de que no la encontró, y quedan \textbf{52 ediciones leídas}. Se leyó en septiembre de 2026 \textbf{sin OCR propio} ---el entorno tampoco traía el modelo español---, con un techo de resolución de \textbf{150 ppi} y una sola versión del texto: la capa del propio Boletín. \textbf{El modo de falla de este año es propio y conviene declararlo}: la capa separa con espacios las letras de palabras enteras ---el 19 por ciento de los tokens del año son de una sola letra, y hay cuarenta y seis hojas donde pasan del 30 por ciento---, de modo que una mención caída en un bloque espaciado se pierde en silencio. El barrido corrió por eso la pasada literal tres veces, sobre el texto tal cual, sobre el texto repegado y sobre el texto sin espacios, y el resultado fue que la letra espaciada no escondió ninguna mención: las tres devuelven las mismas 32 hojas. Se miraron 72 imágenes y se cotejaron 50 columnas de 26 ediciones sin hallar una sola columna traspuesta. \textbf{Lo que sí falló, y falló ocho veces, es la palabra suelta}: el apellido del juez de paz del departamento se imprime en la capa «cacrres» y el topónimo Yacones se pierde del todo, de modo que ni el nombre del único acto firmado en el departamento ni el paraje aparecen en ninguna búsqueda, y los dos aparecieron mirando la imagen. El año dejó \textbf{veintiocho fichas} en 32 hojas, \textbf{ningún acto de dominio}, \textbf{ningún plano} ---diecisiete hojas remiten a planos archivados y ninguna de las diecisiete toca al departamento--- y \textbf{ningún balance municipal}, ni de éste ni de ningún otro municipio. \textbf{Y el noveno es el año 1919}, también agregado al cierre: las ediciones \textbf{738 a 786}, \textbf{49 ediciones} ---los cuarenta y nueve viernes que van del 17 de enero al 19 de diciembre, sin un solo día de salida sin edición dentro de ese tramo y sin ninguna edición doble---. El año tiene \textbf{cincuenta y dos viernes}, y los tres que quedan afuera son los del 3 y el 10 de enero, que caen entre la extraordinaria del 31 de diciembre de 1918 y la primera edición leída, y el del 26 de diciembre, que el repositorio reemplaza por un aviso de que no lo encontró. Con eso, \textbf{440 hojas} sobre los \textbf{478 folios} que la paginación corrida declara. Se leyó en septiembre de 2026, otra vez sin OCR propio del cuerpo, con un techo de 300 ppi. \textbf{Y trae el mejor resultado de cobertura de toda la serie}, por una razón de método que conviene declarar: la cabecera de cada hoja imprime «\textit{Pag.\ N}» y la capa de texto la entrega en once de las 440, pero \textbf{la franja de cabecera se puede reconocer ópticamente con el modelo inglés}, porque un folio de tres cifras no necesita el español, y así se leyó en 307. Con la cadena reconstruida el resultado es tajante: \textbf{dentro de cada una de las 49 ediciones el folio avanza de a uno sin un solo salto}, y las \textbf{38 hojas que faltan} caen todas al final de veintiocho ediciones, que es donde el Boletín de esta época imprime los avisos y los edictos. No falta una sola hoja del medio de ninguna edición. De las nueve ediciones con actos del departamento, cuatro están incompletas. El año dejó \textbf{nueve actos} del departamento, \textbf{ningún acto de dominio}, \textbf{ningún plano} ---diecinueve hojas remiten a uno y ninguna de las diecinueve toca al departamento--- y \textbf{ningún resumen de Tesorería}, contra los tres de 1918. \textbf{Y el décimo es el año 1920}: las ediciones \textbf{788 a 839}, de las que faltan cuatro ---la 793 y la 822, que el repositorio reemplaza por un aviso de que no las encontró, y la 837 y la 838, que directamente no están---, de modo que son \textbf{48 ediciones} y \textbf{526 hojas distintas} sobre seiscientos folios impresos, del 1.081 al 1.680. El techo de resolución vuelve a los \textbf{300 ppi}, el doble que en 1918, y otra vez sin OCR propio del cuerpo. \textbf{La cadena de folios de este año cierra hacia atrás un hueco de 1919}: aquel año termina en el folio 1.072 y 1920 abre en el 1.081, de modo que los ocho folios del medio son la edición \textbf{787}, la del viernes 26 de diciembre de 1919, que ninguno de los dos repositorios tiene. Hay además una \textbf{duplicación del archivo}: la edición 810 trae sus hojas 7 y 8 repetidas como 9 y 10, con el mismo folio impreso y el mismo texto, y por eso las hojas distintas son 526 y no 528. El año dejó \textbf{diez actos} del departamento, en diez hojas y con dieciséis apariciones ---de las diecinueve coincidencias del término, tres son falsos positivos resueltos---, \textbf{ningún acto de dominio} y \textbf{ningún plano} ---catorce hojas remiten a uno y ninguna toca al departamento---; publicó en cambio \textbf{diez resúmenes del movimiento de la Tesorería General}, más que ningún otro de los años cuyos resúmenes se contaron uno por uno ---uno en 1911, cuatro en 1912, tres en 1918 y ninguno en 1919---, y \textbf{ningún balance municipal}. \textbf{Y el undécimo es el año 1921}: las ediciones \textbf{840 a 888}, de las que seis ---la 859, la 872, la 876, la 878, la 884 y la 885--- son marcadores de ausencia del repositorio y no ediciones, de modo que son \textbf{43 ediciones} y \textbf{742 hojas} de Boletín. La cadena de folios corre del \textbf{1.681 al 2.532} y \textbf{empalma sin hueco con el cierre de 1920}, que fue el 1.680: es el único par de años consecutivos de toda la serie cuya paginación se puede seguir sin interrupción. Entre la cadena de folios y las hojas leídas faltan \textbf{110 hojas}: las de las seis ediciones ausentes y ocho más dentro de la edición 852. El reparto entre esas dos causas no se rehízo hoja por hoja y por eso no se declara. Techo de resolución de 283 a 307 ppi, y otra vez sin OCR propio. El año dejó \textbf{dieciocho actos} del departamento, \textbf{dos de ellos de dominio} ---los primeros desde 1917, que el apéndice \ref{cap:dominio} incorpora---, \textbf{ningún plano} y \textbf{ningún balance ni resumen de Tesorería}: después de los diez de 1920, ninguno. \textbf{Y el duodécimo es el primer semestre de 1922}, y es el único tramo de la serie que no cubre un año entero: las ediciones \textbf{890 a 916}, \textbf{27 ediciones} y \textbf{446 hojas} (445 útiles), del 13 de enero al 21 de julio. \textbf{Lo que lo corta no es el método sino el archivo}: de la edición \textbf{917 en adelante no hay PDF} en el repositorio, de modo que cualquier acto del departamento del segundo semestre de 1922 queda fuera de alcance y así se declara. Se leyó en septiembre de 2026 \textbf{sin OCR propio} ---el entorno no traía el modelo español--- y con un \textbf{techo de resolución de 150 a 192 ppi, la mitad que en 1909--1912}, de modo que la única lectura posible fuera de la capa del Boletín fue la imagen: 46 recortes sobre 31 hojas de 16 ediciones, y 31 columnas cotejadas sin hallar una sola columna traspuesta. \textbf{Este tramo tiene tres huecos de especie distinta y conviene separarlos}: la edición \textbf{889} no existe en el repositorio, que la reemplaza por un aviso de que no la encontró; \textbf{no hubo edición el viernes 14 de abril} ---las tapas de la 902, la 903 y la 904, leídas en la imagen, dicen 7, 21 y 28 de abril, y la numeración corre sin salto, de modo que es una semana sin Boletín y no una edición perdida, y el corpus no da la causa---; y \textbf{faltan hojas dentro de las ediciones entregadas}, que es el hueco más incómodo: la cadena de folios impresos ---fijada en la imagen con 890 h.\ 1 = 2.531 y 890 h.\ 18 = 2.548--- avanza más rápido que el conteo de hojas de los PDF, con un desvío acumulativo y del mismo signo en los cinco tramos medidos, que da \textbf{alrededor de trece hojas de más} entre las ediciones 890 y 910. En qué ediciones concretas faltan no se estableció. \textbf{El modo de falla de la capa es la palabra suelta}, otra vez: cinco casos medidos, cuatro de ellos decisivos para una ficha ---«galdefa» y «Cald\#ra» por Caldera, «c aldera» partida en dos, «Masc-ief» por Masclef y el nombre de la finca del remate, que la capa devuelve como una cadena de signos y la imagen dice «Potrero del Gallinato»---. \textbf{Y por tercer año consecutivo, la unión correcta del corte de renglón rescata un acto del departamento}: la hoja 10 de la edición 899, que trae el decreto 349 sobre la comisaría de policía, no aparece en ninguna búsqueda literal si el guion blando se quita como carácter suelto en lugar de quitarse junto con el salto de renglón. El tramo dejó \textbf{veintitrés actos} del departamento en veinticuatro hojas ---veintidós por el nombre del departamento y uno, el decreto 176 sobre el subcomisario de Mojotoro, sólo por el nombre del paraje---, \textbf{ningún acto de dominio} salvo el aviso de remate del Potrero del Gallinato, \textbf{ningún plano} ---tres remisiones a planos en todo el corpus y ninguna del departamento, contra las dos de 1912--- y \textbf{ningún balance municipal}, ni de éste ni de ningún otro municipio. \textbf{Publicó tres resúmenes del movimiento de la Tesorería General} ---abril, mayo y junio---, cuya cadena de saldos cierra al centavo entre los tres, \textbf{con el rótulo del tercero equivocado}: encabeza «A SALDO DE JULIO DE 1922» una cifra que es la de mayo, el mismo accidente que 1912, y la cadena se arma por la cifra y no por el rótulo. Ninguno de los tres desagrega por departamento. \textbf{Y el decimotercero es el año 1923 entero}: las \textbf{52 ediciones} de la \textbf{939 a la 990} y \textbf{673 hojas}, una edición por cada uno de los cincuenta y dos viernes del año, \textbf{sin un solo hueco en la serie numérica, sin ningún día de salida sin edición y sin ninguna edición doble}. Era, en ese sentido, el tramo más completo de la serie hasta que se incorporó 1924. Se leyó en septiembre de 2026 \textbf{sin OCR propio} ---el entorno no traía el modelo español--- y con un \textbf{techo de resolución de 150 ppi}, uniforme en todo el año: cada hoja trae dos imágenes incrustadas sobre el mismo rectángulo, una de 150 ppi y otra de 75, de modo que pedir más resolución agranda píxeles y no agrega información. Las \textbf{52 tapas se leyeron en la imagen} ---las 32 primeras están compuestas en tipo de fantasía que la capa no resuelve--- y cincuenta y una coinciden con el viernes que les toca; la excepción es la \textbf{edición 945}, cuya tapa dice «\textit{Viernes 17 de Febrero de 1923}» cuando el 17 fue sábado y el viernes de esa semana fue el 16: el dígito se verificó contra otras dos tapas del mismo año y \textbf{no se concilia}. \textbf{El control de folios existe hasta la edición 970 y después desaparece}: la cadena no reinicia en enero ---la edición 939 abre en el folio 3.275, heredado de 1922, al revés de lo medido en 1911 y 1912--- y corre corrida hasta el 10 de agosto, donde marca \textbf{catorce folios de más} sobre las hojas presentes, repartidos en nueve tramos; \textbf{quince hojas del año están foliadas y no están en el repositorio}, y la única que puede nombrarse es el folio 3.334. \textbf{Desde la edición 971 el folio reinicia en 1 en cada edición}, de modo que \textbf{322 de las 673 hojas quedan sin ese control} y no hay manera barata de saber si esas veinte ediciones están completas. El año editorial cambia dentro del corpus: ANO XIV en las ediciones 939 y 940 y ANO XV desde la 941. \textbf{El modo de falla de este año es propio y es el más grave de la serie: el original imprime el nombre del departamento con las letras separadas}, «C A L D E R A» y «C ald era», de modo que ninguna búsqueda literal lo devuelve; la pasada laxa, reescrita para admitir espacios entre todas las letras, \textbf{rescató tres actos que la literal no veía}, entre ellos el único aviso de dominio del año. La pasada difusa, que en 1909, 1911 y 1912 rescataba menciones, \textbf{en 1923 no rescata ninguna}: sus seis formas son todas falsos positivos. El año dejó \textbf{diecisiete actos} del departamento en veintiuna hojas, más uno que no lo nombra y que entró por el anillo de topónimos ---el cierre del camino a la Quebrada de Lesser---; \textbf{ningún plano} ---las 673 hojas son texto a dos columnas, sin una sola lámina---, \textbf{ningún balance municipal} y ninguna cifra propia del departamento fuera del sueldo de \$44 mensuales del encargado del Registro Civil, que el presupuesto fija por igual para más de cuarenta oficinas de campaña y no desglosa. \textbf{Y el decimocuarto es el año 1924 entero}: las \textbf{52 ediciones} de la \textbf{991 a la 1042} y \textbf{792 hojas}, una por cada viernes del año, y \textbf{es el primer año de toda la serie barrida en que el calendario cierra solo}: cincuenta y dos viernes, cincuenta y dos ediciones, sin huecos, sin ediciones dobles y sin una sola tapa cuya fecha no caiga donde debe ---sólo tres ediciones imprimen un día de la semana que no corresponde a su propia fecha, y se declaran---. Se leyó en septiembre de 2026 \textbf{sin OCR propio} y con el mismo techo de \textbf{150 ppi} de 1923. \textbf{El control de completitud cambia de naturaleza y mejora}: en 1924 la foliatura no es corrida en el año ---reinicia en cada edición--- pero cada hoja imprime su número de página dentro de la edición y \textbf{el sumario de la hoja 1 remite a esas páginas}, de modo que la página más alta citada por el sumario nunca supera las hojas presentes: en las 52 ediciones el cotejo cierra. El año editorial cambia entre el 11 y el 18 de enero, en la tercera edición. El tramo dejó \textbf{veinticuatro actos} del departamento en veinte ediciones ---tres menos que 1918 y uno menos que 1925---, \textbf{ningún plano} y \textbf{ningún balance municipal}, ni de éste ni de ningún otro departamento. \textbf{Publicó once resúmenes mensuales de la Tesorería General} y el \textbf{Presupuesto General} entero, que ocupa once hojas: se recalcularon los doce cuadros sobre la imagen y \textbf{once cierran exactos}, incluidos los sub-subtotales internos. \textbf{La excepción es marzo}, cuyo renglonaje de egresos suma \$398.121,34 contra un subtotal impreso de \$382.121,34: \textbf{una diferencia de exactamente dieciséis mil pesos}, y el error está en el detalle y no en el subtotal, porque el subtotal impreso es el que cierra con el saldo y con la columna de ingresos. \textbf{Y el decimoquinto es el año 1925 entero}: las \textbf{52 ediciones} de la \textbf{1043 a la 1094} y \textbf{798 hojas}, una por cada viernes y sin huecos en la serie, como 1924. Mismo techo de \textbf{150 ppi} y \textbf{sin OCR propio}. \textbf{El control de completitud es el mismo que en 1924 y su alcance conviene declararlo con precisión}: la paginación reinicia en cada edición y la cabecera imprime «Pág.\ N»; \textbf{la capa alcanza a leer ese número en 24 de las 746 hojas interiores del año, y en las 24 coincide con el número de hoja dentro de su edición, sin una sola discrepancia}. Sobre esa base no hay indicio de hoja interior faltante, y el resto no es ilegible en el papel sino en la capa, porque el reconocimiento destroza el renglón de cabecera: \textbf{es una comprobación sobre el 3\,\% de las hojas y no sobre todas}, y así se usa. El tramo dejó \textbf{veinticinco actos} del departamento, \textbf{ningún plano} y \textbf{ningún balance municipal}; publicó nueve cuadros con subtotal, todos recalculados. \textbf{Una sola ficha del año no tiene lectura de imagen} ---el edicto y el auto de concesión del cateo 1168-C, leídos en la capa--- y se declara. \textbf{Y el decimosexto es el año 1926 entero y el decimoséptimo es 1927}: las ediciones \textbf{1096 a 1147} con \textbf{769 hojas} y las \textbf{1148 a 1199} con \textbf{745}, las dos series semanales completas, una edición por cada uno de los cincuenta y dos viernes de cada año y sin un solo hueco. Mismo techo de \textbf{150 ppi} y sin OCR propio. \textbf{Los dos años trajeron una comprobación que los anteriores no permitían}: en 1926 el nombre del gobernador se leyó en la imagen porque la capa lo devuelve como «CORBACHO» y «CORBfltón», y en 1927 la declaración de los días de salida \textbf{no está en la tapa sino en la tarifa del propio Boletín}, que dice «aparece los Viernes» ---quien la busque en la tapa, como en 1911 y 1912, no la encuentra---. \textbf{Las cadenas de saldos de los dos años son incompletas y se declaran así}: 1926 publica cinco resúmenes de Tesorería que no son meses consecutivos y con los que \textbf{no se puede armar ninguna cadena} ---uno de ellos, en abril, rotula «noviembre ppdo.», que es imposible---; 1927 publica nueve, los nueve recalculados y los nueve cerrados, pero \textbf{sin junio, julio, agosto ni diciembre}, de modo que la cadena del año se arma en dos tramos con un puente de una sola cifra, el saldo de agosto que el resumen de septiembre arrastra. \textbf{Y 1927 dejó la lección de método del año}: el remate de la finca Los Sauces sale \textbf{tres veces con tres números de aviso distintos}, de modo que cotejar números no lo habría reunido y cotejar texto sí. \textbf{Y el decimoctavo es el año 1928 entero}: las ediciones \textbf{1200 a 1251}, una por cada viernes, con el mismo techo de \textbf{150 ppi} y sin OCR propio. Dejó \textbf{veintitrés actos} del departamento. \textbf{Su cadena de saldos tiene dos huecos declarados} ---faltan marzo y octubre--- y \textbf{varios resúmenes no imprimen renglón de cierre}, de modo que la cadena se arma por la cifra de apertura del mes siguiente y no por el cierre del anterior. \textbf{Y el año trae la anomalía de numeración más grande de toda la serie}: en la misma hoja y con la misma fecha, la serie de decretos pasa del \textbf{7097 al 8001}, novecientos números en un día, con los dos tramos firmados por el mismo gobernador y el mismo ministro. \textbf{El barrido lo declara y no lo concilia}; el capítulo \ref{cap:politica} propone una lectura ---que la cifra de los miles avanza cada cien decretos--- que explica el patrón y no está probada. Las dos quedan escritas. \textbf{Y el decimonoveno es 1929, el vigésimo es 1930, el vigésimo primero es 1931 y el vigésimo segundo es 1932, los cuatro enteros}: las ediciones \textbf{1252 a 1303}, \textbf{1304 a 1355}, \textbf{1357 a 1407} y \textbf{1408 a 1461}, con el mismo techo de \textbf{150 ppi} y sin OCR propio en ninguno de los cuatro ---el entorno no trae el modelo español, de modo que todo lo que la capa no resuelve se lee mirando la imagen y no reconociéndola---. \textbf{Los cuatro se bajaron del Release del repositorio \texttt{boletines-salta} y ninguno trajo salida de barrido automático}: no hubo pre-informe, ni CSV, ni \texttt{EXTRAIDO.json}, de modo que \textbf{se trabajó con una sola versión del texto} y el cotejo de completitud se hizo contra la paginación impresa, que en esta época reinicia en cada edición. \textbf{Y hay un desajuste de insumos que conviene dejar dicho porque afecta lo que el barrido ve pasar}: el \texttt{lee\_config.json} publicado en esos Releases está atrasado ---cierra en 1915, con 180 actores y 70 términos de anillo 2--- mientras que el vigente llega a 1928, con 388 actores y 139 términos; los barridos de 1930 a 1932 usaron el vigente y lo declararon, pero los doscientos ocho actores de 1916 a 1928 que el Release no tiene son justamente los de la época más cercana a esos años. \textbf{Ninguno de los cuatro años publica un solo plano}: las 1.321 hojas de 1931 son texto a dos columnas sin una lámina, y las sesenta y nueve hojas que mencionan la palabra «plano» remiten siempre a uno archivado en un expediente. \textbf{Y el trigésimo tercero es 1947 y el trigésimo cuarto es 1948}, los dos agregados al cierre y barridos en septiembre de 2026: \textbf{158 archivos y 2.585 hojas} el primero, de la edición \textbf{2869 a la 3026}; \textbf{248 archivos y 4.674 hojas} el segundo, desde la \textbf{3060}, que es el tramo más voluminoso de la serie temprana. \textbf{Los dos están en un estado distinto del resto y conviene declararlo con precisión}: de ellos hay pre-informe automático y \textbf{no hay lectura sobre la imagen}. La confianza media del reconocimiento propio es de \textbf{62,0} en 1947 y de \textbf{66,5} en 1948; \textbf{ninguna hoja de ninguno de los dos años alcanza el umbral de publicable} ---0 publicables sobre 2.585 hojas en 1947 y 0 sobre 4.674 en 1948---, y quedan marcadas para mirar enteras \textbf{90 hojas} de 1947 y \textbf{187} de 1948. El lector automático de tapas de 1947 devuelve números de edición falsos ---«Nº 1805», «Nº 4780», «Nº 28911»--- en buena parte del año, y una tapa imprime «4947». \textbf{De modo que todo lo que el capítulo \ref{cap:siglo} consigna de 1947 y 1948 es lectura de capa de texto y está sujeto a verificación sobre el facsímil}, y así va declarado ahí y pedido en el apéndice \ref{cap:pedidos}. Mientras eso no se haga, \textbf{estos dos años no cuentan para la regla de que dentro de la ventana un cero es un resultado}: cuentan sus hallazgos, no sus ausencias. \textbf{Y el trigésimo quinto es 1949}, agregado después: \textbf{281 archivos y 4.869 hojas}, el tramo más voluminoso de toda la serie temprana, con la misma reserva que los dos anteriores. Sus tapas repiten la patología de 1947 ---«Nº 4780» tres veces, «Nº 1805», y fechas imposibles como «1349-01-07» y «3949-01-20»---, de modo que \textbf{el número de edición se toma del nombre del archivo y no de la tapa}. \textbf{Y el trigésimo sexto es 1950}: \textbf{277 archivos y 4.156 hojas}, con la misma reserva. \textbf{Y el vigésimo tercero es 1937, el vigésimo cuarto es 1938, el vigésimo quinto es 1939, el vigésimo sexto es 1940, el vigésimo séptimo es 1941, el vigésimo octavo es 1942, el vigésimo noveno es 1943, el trigésimo es 1944, el trigésimo primero es 1945 y el trigésimo segundo es 1946}, y traen un cambio material que conviene declarar antes que nada: \textbf{los PDF de esta época ya vienen con capa de texto del proveedor, y es buena}. Medido en 1937: de \textbf{1.917 hojas en 53 ediciones}, \textbf{1.897 (el 99\,\%) tenían capa útil} ---al menos cuatrocientos caracteres y más de la mitad de letras--- y sólo veinte necesitaron reconocimiento propio. Las tres pasadas coinciden entre sí y los contextos se leen limpios, sin las deformaciones que dominaron de 1909 a 1932. \textbf{Eso tiene dos consecuencias y las dos se declaran.} La primera: el año entero se barrió en \textbf{dos minutos y medio} contra las horas que costaba un año de la época anterior. La segunda, que es un \textbf{límite}: de las hojas con capa útil no se guardó imagen, de modo que \textbf{este tramo no tiene confianza por hoja, ni medición de tinta, ni deduplicación por \textit{dhash}}; mirar una hoja sigue siendo posible y barato, porque se rasteriza desde el PDF en el momento. \textbf{Y el lector automático de tapas no funciona con esta capa}: devuelve números de edición falsos ---«Nº 2927», «Nº 81871»--- y no resuelve la fecha en la mitad de los casos, de modo que las tapas se leyeron en la imagen, como siempre.

\textbf{Un cero que sólo se mide sumando los tramos}: en los siete años que van de 1910 a 1916, leídos todos edición por edición, el Boletín \textbf{no publica un solo balance municipal}, ni de este departamento ni de ningún otro municipio de la provincia. Cada barrido lo declara para su año; sumados, son siete años consecutivos (capítulo \ref{cap:siglo}). \textbf{El barrido de 1918 vuelve a declarar el mismo cero}, y el año publica en cambio tres resúmenes mensuales del movimiento de la Tesorería General ---septiembre, octubre y noviembre--- que encadenan sus existencias sin un centavo de diferencia: el Estado provincial rinde su caja mes a mes y ningún municipio rinde la suya. El resto de ese hueco dejó de ser parcial al cierre: \textbf{1919, 1920 y 1921, el primer semestre de 1922 y 1923 entero se leyeron después edición por edición}, y 1920, 1921, 1922 y 1923 declaran el mismo cero cada uno por su cuenta. En los tramos leídos se midió el desfase entre la
página impresa y la del escaneo, se reconstruyeron los asuntos cortados entre páginas, se
cotejó la capa de texto contra la imagen donde el OCR falló y se descartaron los falsos
positivos de uno en uno.

La diferencia práctica es ésta: \textbf{en 1908--1921 y en 1923--1946, y sólo ahí, este libro
puede afirmar que algo no aparece} ---con la salvedad de los números y las páginas faltantes---, y también en el \textbf{primer semestre de 1922}, que se leyó del mismo modo y queda entre las dos ventanas. En el \textbf{segundo semestre de 1922} ---del que el repositorio no tiene un solo PDF--- y en todo lo posterior a 1946, la ausencia sigue significando solamente
que el barrido no lo encontró. Dentro de esas dos ventanas, en cambio, un cero es un resultado
medido y no un hueco de búsqueda. El río Trancas, por ejemplo, no aparece en ninguna de las
ediciones leídas ---las únicas «Trancas» del período son un paraje de Orán y la localidad
tucumana---, y sólo entra al archivo en 2015. Y las cincuenta ediciones que van del
B.O.\ Nº 2.115 al Nº 2.164 ---del 5 de octubre al 5 de diciembre de 1944, 441 páginas--- no
nombran una sola vez a La Caldera ni a La Calderilla: la única mención del departamento en esos
dos meses es la Colonia de Vaqueros, en el Decreto 4947-G del 18 de octubre (B.O.\ Nº 2.126).
La ventana se cierra dos ediciones antes de la que publica el Revalúo General con la nómina del departamento (B.O.\ Nº 2.166).

\section{Expedientes judiciales}

\begin{itemize}[leftmargin=*]
\item \textbf{Amparo colectivo}, Expte.\ Nº MIN 23860/19 ---otras piezas lo numeran EMI 23860/2020; la discrepancia se declara en el capítulo \ref{cap:metodo}---, CUIJ 17-00023860-0, caratulado
«Lazarte Vigabriel vs.\ Municipalidad de La Caldera y Ministerio de Producción, Trabajo y
Desarrollo Sustentable de la Provincia de Salta». Analizado en el capítulo \ref{cap:amparo}.
\item \textbf{Incidente de cumplimiento}, Expte.\ Nº INC 23860/1, con las piezas que las
demandadas acompañaron en cumplimiento de la condena. De este incidente provienen siete de las treinta y ocho láminas de este volumen.
\item \textbf{Ejecución de sentencia}, Expte.\ Nº 780280.
\item \textbf{Documentos de 2022--2023 del expediente}: nota de la Secretaría de Recursos Hídricos del 20/01/2023 (expte.\ 0010007-67527/2019-19) con el Plan de Monitoreo de Aguas Superficiales 2022; informe y plano de Aguas del Norte (02/12/2022); diseño del taller participativo (31/08/2023); y presentación del Plan Maestro de Drenaje Pluvial--Fluvial, sin fecha y sin expediente citado ---la Advertencia preliminar la cuenta entre los documentos del Plan de Manejo y no entre las piezas del expediente--- (láminas \ref{lam:pmdp} y \ref{lam:pmdpcotas}).
\end{itemize}

\section{Expedientes administrativos}

Se citan en el cuerpo con su número completo y se reúnen, por organismo, en el apéndice
\ref{cap:pedidos}. Los de consulta más extensa son el \textbf{0090136-221704/2023-2} ---
Plan Integral de Manejo de la Microcuenca del Río La Caldera, emitido en octubre de 2024 ---
y los expedientes de determinación de línea de ribera del período 2013--2015, en particular
el \textbf{34-167.127/13}.

\begin{cautela}
\textbf{Ningún expediente se consultó físicamente.} Todo lo que este libro sabe de ellos
proviene de lo publicado en el Boletín, de lo incorporado a piezas judiciales, o de lo que
otros documentos transcriben. Es el límite más serio del relevamiento y está declarado
también en el capítulo \ref{cap:metodo}.
\end{cautela}

\section{Cartografía}

\begin{itemize}[leftmargin=*]
\item \textbf{«Informe de mapas de cuenca de río La Caldera»}, Secretaría de Recursos
Hídricos de Salta, mayo de 2022, ocho hojas con tres mapas temáticos a escalas 1:125.000 y
1:40.000 en POSGAR 2007 --- Faja 3. Reproducido parcialmente en las láminas
\ref{lam:mapa1}, \ref{lam:mapa2}, \ref{lam:mapa3} y \ref{lam:mojotoro}.
\item \textbf{Atlas del Código de Ordenamiento Territorial} del Municipio de La Caldera,
ocho planos, publicados sólo como imágenes dentro del Código y sin sus capas. Dos se reproducen en las láminas \ref{lam:anexo1} (tomada del expediente judicial) y \ref{lam:anexo2} (tomada del PDF municipal).
\item \textbf{Visor de IDESA}, consultado el 15 y el 16 de septiembre de 2026: capas de límites municipales y de catastro, en las láminas \ref{lam:idesalim}, \ref{lam:idesamun}, \ref{lam:idesacat} y \ref{lam:cintas}.
\item \textbf{Plano de mensura 000962} del departamento, en la lámina \ref{lam:plano962} y su detalle vial en la lámina \ref{lam:zonacamino}.
\item \textbf{Capas del informe técnico del OTBN (Ley 8483)}: capa de cuencas con parámetros, ráster del mapa final y catastro rural, descargadas del repositorio que cita la nota 7 del informe; y \textbf{cartografía de radios del Censo 2022 (INDEC)}, cuyos veintiocho radios forman el límite departamental usado. Con ellas y con las coordenadas del anexo de la Resolución 023/14 se elaboró la lámina \ref{fig:otbn}; el resultado se verificó con el \textbf{límite departamental del IGN (SIG 250)}, en la versión GeoJSON del repositorio público \textit{departamentos\_argentina}.
\item Capas de base del \textbf{Instituto Geográfico Nacional} y del geoportal
\textbf{IDESA} de la provincia de Salta, y modelo de elevación \textbf{SRTM}, según declara
el propio informe de 2022.
\end{itemize}

\section{Cómo se produjo cada lámina propia}

Diecisiete de las treinta y ocho láminas son de elaboración propia: \textbf{diez mapas}
---\ref{fig:ubicacion}, \ref{fig:fincas}, \ref{fig:parajesmapa}, \ref{fig:riberaedictos},
\ref{fig:catastro}, \ref{fig:planos}, \ref{fig:formacionradios}, \ref{fig:kondorwaira},
\ref{fig:otbn} y \ref{fig:agua}--- y \textbf{siete gráficos} ---\ref{fig:parajesnom}, \ref{fig:caudales},
\ref{fig:formacioncomparacion}, \ref{fig:padron}, \ref{fig:cargos}, \ref{fig:coparticipacion} y \ref{fig:defensas}---. Todas salen del \textbf{repositorio cartográfico \texttt{mapas\_caldera}}
---\texttt{github.com/ediedrich/mapas\_caldera}---, que acompaña a este libro.

\textbf{Conviene decir con exactitud qué contiene y qué no}, porque la diferencia decide
quién puede rehacer una lámina. El repositorio versiona \textbf{los scripts}, \textbf{las
figuras resultantes} y \textbf{las capas propias}: las tres tablas de cargos y la planilla de
ribera, que son transcripción de este trabajo y de nadie más. \textbf{Las capas de terceros
---catastro de IDESA, radios del INDEC, límites, hidrografía, red vial, relieve y
asentamientos del IGN y de BAHRA, y el ordenamiento de bosques--- no se redistribuyen}: el
repositorio declara de cada una su origen y la fecha en que se consultó, y trae un script
---\texttt{scripts/bajar\_capas.py}--- que las baja de un \textit{Release} aparte. Es una
decisión deliberada y de licencia: el archivo envejece, la procedencia no.

\textbf{Ese \textit{Release} contiene treinta y nueve capas} ---catastro, radios censales,
límites, hidrografía, red vial, relieve, asentamientos y las tres del ordenamiento de
bosques---, todas recortadas al rectángulo del departamento y reproyectadas a EPSG:4326, y
pesa \textbf{47,9 MB}. El repositorio trae además un \texttt{requirements.txt} con las
versiones de biblioteca con que se probó.

\textbf{Y la cadena se ejecutó entera antes de cerrar este volumen}, que es la única forma de
afirmar que funciona: en una máquina donde ninguna de esas capas existía, se clonó el
repositorio, se corrió \texttt{bajar\_capas.py} y se corrieron los nueve scripts. \textbf{Los
nueve terminaron y las nueve láminas se rehicieron.} Las que están publicadas en este libro
son ésas.

\textbf{Qué script produce cada una}:

\begin{itemize}[leftmargin=*,nosep]
\item \texttt{scripts/mapa\_ubicacion.py} --- lámina \ref{fig:ubicacion}.
\item \texttt{scripts/mapa\_fincas.py} --- lámina \ref{fig:fincas}.
\item \texttt{scripts/tabla\_parajes.py} --- lámina \ref{fig:parajesnom}; la tabla de doce
nóminas que arma este script es además \textbf{la fuente única} de la lámina siguiente.
\item \texttt{scripts/mapa\_parajes.py} --- lámina \ref{fig:parajesmapa}.
\item \texttt{scripts/mapa\_ribera.py} --- lámina \ref{fig:riberaedictos}, sobre la planilla
\texttt{datos/ribera/} que transcribe cada punto con su número de Boletín, su orden de
publicación y su fecha de consulta. \textbf{Sólo esa lámina}: hasta la revisión de este
apéndice, estas líneas atribuían a este script también la lámina \ref{fig:catastro}, y no es
cierto ---el script tiene una sola salida---. La \ref{fig:catastro} usa la misma planilla y el
mismo estilo, pero \textbf{su script no está publicado}, y por eso pasa a la lista de abajo.
\item \texttt{scripts/mapa\_planos.py} --- lámina \ref{fig:planos}.
\item \texttt{scripts/mapa\_kondorwaira.py} --- lámina \ref{fig:kondorwaira}.
\item \texttt{scripts/mapa\_agua.py} --- lámina \ref{fig:agua}, sobre las capas de cursos y límites del IGN y el campo «finca» del catastro; los dos paneles y la línea de tiempo salen de la tabla de actos que el propio script declara.
\item \texttt{scripts/grafico\_cargos.py} --- lámina \ref{fig:cargos}, sobre las tres tablas de
\texttt{datos/cargos/}: \texttt{titulares.csv}, \texttt{comision.csv} y \texttt{tutela.csv}.
\textbf{Dos columnas de esas tablas}, \texttt{exacta\_desde} y \texttt{exacta\_hasta}, dicen si
cada borde de segmento sale de un acto publicado o si el archivo no lo fecha, y es lo que el
dibujo traduce en bordes a tope o degradados; el script informa la proporción al correr.
\item \texttt{scripts/base.py} --- estilo común, contorno, escala y norte de todas ellas.
\item \texttt{scripts/grafico\_coparticipacion.py}
--- lámina \ref{fig:coparticipacion}, sobre
\texttt{datos/fiscal/}\allowbreak\texttt{coparticipacion\_1947.csv}, que transcribe la tabla del
decreto 5276-E municipio por municipio, con su recaudación de 1946 y su porcentaje.
\item \texttt{scripts/grafico\_defensas.py} --- lámina \ref{fig:defensas}, sobre \texttt{datos/fiscal/defensas\_1947.csv}, que transcribe los nueve ríos del Ítem 6 de la Ley 834 con su monto.
\item \textbf{Y seis láminas propias no tienen todavía su script en el repositorio}: el de las líneas de ribera sobre el catastro (\ref{fig:catastro}), el mapa
del ordenamiento de bosques (\ref{fig:otbn}), el de formación por radio censal
(\ref{fig:formacionradios}) y los tres gráficos de caudales, formación comparada y padrón
(\ref{fig:caudales}, \ref{fig:formacioncomparacion} y \ref{fig:padron}). Las fuentes
de las que salen están declaradas en cada epígrafe, de modo que son verificables; \textbf{pero
no son reproducibles paso a paso hasta que el script que las dibuja se incorpore}. Se declara
aquí porque es exactamente la diferencia que este apéndice acaba de reclamar para las otras
once.
\end{itemize}

\textbf{Todas las capas están en EPSG:4326} y las superficies se calculan en proyección
equivalente (EPSG:6933); cada capa nacional se recortó al rectángulo 65,85º--65,0º O,
24,85º--24,25º S, que contiene al departamento con margen, con el script
\texttt{scripts/incorporar.py}, que deja constancia de cuántos elementos quedaron dentro. El
\texttt{LEEME} del repositorio trae el inventario capa por capa con su origen y su fecha de
consulta.

\textbf{De modo que el estado de las diecisiete láminas propias, al cierre de este volumen, es
éste}: las diecisiete son \textbf{verificables} ---cada epígrafe declara con qué datos se hizo, y
esos datos tienen origen y fecha---; \textbf{once son además reproducibles} ---su script
está publicado, sus capas se bajan del \textit{Release} y la cadena se probó completa---; y
\textbf{seis no lo son todavía}, porque su script no está. \textbf{Verificable no es lo
mismo que reproducible}, y este apéndice prefiere distinguirlas antes que dejar que el lector
las confunda.

\subsection*{La imagen de cubierta, que no es una lámina y se declara igual}

La imagen que abre el libro no integra el índice de láminas, no ilustra ningún argumento y no
se cita en ningún capítulo. Se declara acá porque una imagen sin procedencia declarada es, para
el lector, indistinguible de un documento.

Es \textbf{de autoría del autor de este libro} y su cadena de producción es ésta: dibujos a mano,
digitalización, composición y retoque en \textsc{corel draw}, y una pasada final de
\textbf{optimización con Gemini}, un modelo generativo de imagen. \textbf{No es una fotografía,
no documenta un lugar determinado y no debe leerse como registro de nada}: es una composición
---un valle con un trazado de amanzanamiento superpuesto--- que enuncia el objeto del libro y no
lo prueba. Ninguna de las treinta y ocho láminas del volumen pasó por un modelo generativo, ni
en su producción ni en su retoque.

\section{Series estadísticas}

\begin{itemize}[leftmargin=*]
\item \textbf{INDEC}, Censo Nacional de Población, Hogares y Viviendas, ediciones 2010 y
2022. Procesamiento propio con \textbf{Redatam} sobre el área 66077 y sus fracciones.
\item \textbf{Censo Nacional Agropecuario}, ediciones 1988, 2002, 2008 y 2018, para la superficie
implantada y la cantidad de explotaciones del departamento.
\item \textit{Estadística Hidrológica de la República Argentina}, edición 2004, y
\textbf{Atlas Digital de los Recursos Hídricos Superficiales} del Instituto Nacional del
Agua, citados en el Plan de Manejo como fuente de las series de aforo de las cinco estaciones de la cuenca.
\item \textbf{ENARGAS}: atlas provincial del gas, lámina de Salta, sin fecha; mapa nacional de localidades abastecidas con gas por redes, 10/2024; y exportación del mapa interactivo, 16/09/2026. Reproducidos en las láminas \ref{lam:gasatlas}, \ref{lam:gasnacional} y \ref{lam:gasvisor} (capítulo \ref{cap:redes}).
\item \textbf{Censo 2022 por radio censal}: nivel más alto que cursó, completitud de ese nivel y edad simple, para los veintiocho radios del departamento, procesados en Redatam en septiembre de 2026; y la \textbf{cartografía de radios censales 2022} en su versión corregida por Gonzalo M.\ Rodríguez (CEUR-CONICET, versión 1.0), que ubica cada radio (capítulo \ref{cap:poblacion}, lámina \ref{fig:formacionradios}).
\item \textbf{Tribunal Electoral de la Provincia de Salta}: escrutinio definitivo de las elecciones generales del 11 de mayo de 2025, planilla por mesa y por lista, categoría senador, departamento La Caldera; planilla de establecimientos y mesas de 2025, con los empadronados por mesa; y los listados de electores por circuito de 2017, 2019, 2021 y 2023 (capítulo \ref{cap:politica}).
\item \textbf{Sistema Nacional de Información Hídrica} (SNIH), Subsecretaría de Recursos Hídricos de la Nación, Red Hidrológica Nacional: datos históricos de caudal medio mensual de las estaciones 623 (El Volcán), 625 (El Angosto), 634 (San Alejo), 635 (Santa Rufina) y 640 (Desembocadura al Nieves), 1942--1986, descargados en septiembre de 2026 (capítulo \ref{cap:ribera}). La exportación trae algunos registros con una hora de desfase ---fechados a las 23 del último día del mes anterior---, que este libro corrigió antes de calcular.
\end{itemize}

\begin{cautela}
La \textit{Estadística Hidrológica} y el \textit{Atlas Digital} \textbf{no se consultaron en su edición original}: este libro las conoce por la cita que de ellas hace el Plan de Manejo. Las series de caudal medio mensual, en cambio, se descargaron del SNIH y son de primera mano; lo que sigue dependiendo de aquellas compilaciones son los caudales máximos instantáneos, que el SNIH no ofrece.
\end{cautela}

\section{Instrumentos normativos de referencia}

\begin{itemize}[leftmargin=*]
\item \textbf{Ley provincial 8483}, publicada el 3 de enero de 2025, revisión del
Ordenamiento Territorial de Bosques Nativos, y el informe final del equipo técnico que
constituye su Anexo I.
\item \textbf{Ley nacional 26.160} y las resoluciones de relevamiento territorial del
Instituto Nacional de Asuntos Indígenas correspondientes a la provincia de Salta.
\item \textbf{Código de Aguas de la Provincia de Salta}, Ley 7017, y su reglamentación.
\end{itemize}

\section{Partes de prensa oficiales}

El sitio de la Secretaría de Comunicación de la Provincia publica partes de prensa fechados y
atribuidos, que este libro usa en varios lugares. \textbf{No son fuente documental}: acreditan
que una reunión o un anuncio existieron y qué se dijo, no que lo anunciado se haya hecho.
Cada vez que se los cita, la caja de cautela correspondiente lo aclara.

El más usado es el del \textbf{16 de agosto de 2017}, sobre la coordinación de la
determinación de línea de ribera para el Acueducto Norte, en el capítulo \ref{cap:ribera}.

A ellos se agrega el \textbf{sitio de la Municipalidad de La Caldera}, del que proviene la
nota del \textbf{22 de octubre de 2020} sobre las obras en los arroyos El Durazno y
Guaranguay, en el capítulo \ref{cap:amparo}. Vale la misma advertencia, y una más: es la
\textbf{única descripción concreta} que este relevamiento tiene de un trabajo de
encauzamiento realizado por el municipio, y proviene de la comunicación institucional del
propio municipio que lo hizo.

\begin{cautela}
Vale una nota sobre el Boletín Oficial que conviene tener presente para cualquier
continuación. \textbf{Cada instrumento publicado tiene, además del facsímil, una versión en
texto} encabezada por la leyenda «este texto es copia fiel del Boletín Oficial». Eso significa
que los edictos y resoluciones de las ediciones digitalizadas \textbf{son buscables sin
reconocimiento óptico}, y que una afirmación de ausencia sobre ese tramo tiene mucho más
respaldo que sobre los años que solo existen como imagen escaneada.
\end{cautela}

\section{El Boletín Oficial Municipal de la ciudad de Salta}

Este libro usa, en el capítulo \ref{cap:infraestructura}, varias piezas publicadas en el \textbf{Boletín Oficial Municipal de la ciudad de Salta}, que es un registro distinto del Boletín Oficial de la Provincia y del municipal de La Caldera.

\textbf{Se consultaron pocas ediciones}, entre ellas la 2.646 del 2 de octubre de 2024 y la 2.725 del 27 de agosto de 2025, cada una por una razón puntual; la de la segunda fue verificar un rol de una empresa que el libro ya documentaba en otros dos registros. \textbf{No se hizo un relevamiento de esa fuente}, que publica con firma digital desde 2016 y cuyo contenido sobre contrataciones de obra podría ser relevante para este trabajo.

\section{Medidas tomadas sobre imagen}

El capítulo \ref{cap:loteo} contiene un apartado ---el de las parcelas en cinta--- cuyas magnitudes no salen de ningún instrumento sino de \textbf{medir sobre la grilla kilométrica de una exportación del visor provincial}. Largos, anchos y superficie son aproximados, con un margen que este libro no cuantifica y que no es despreciable.

\textbf{Es la única parte del libro construida así}, está señalada como tentativa en su propio encabezado, y el apéndice \ref{cap:pedidos} contiene el pedido que la confirma o la elimina. \textbf{Mientras ese pedido no vuelva respondido, ese apartado no es un hallazgo de este libro.}

\section{Cartografía censal}

La identificación de las tres fracciones censales del departamento ---capítulo \ref{cap:ausencias}--- se hizo sobre la \textbf{capa vectorial de fracciones censales del organismo censal}, obtenida por su servicio público de datos geográficos, y se verificó por dos vías independientes: \textbf{punto en polígono} sobre coordenadas ya establecidas en este libro, y \textbf{coincidencia aritmética} entre la población de cada fracción y la de cada municipio.

\textbf{Es una de las operaciones geométricas del libro} ---las otras son la medición sobre grilla del capítulo \ref{cap:loteo}, el cálculo de distancias entre coordenadas y el cruce de capas del ordenamiento de bosques del capítulo \ref{cap:opacidad}, la más elaborada---, y es reproducible: la capa es descargable y el cálculo, elemental.

\section{Procesamiento censal propio}

Se usaron además dos capas geográficas: la de \textbf{radios del Censo 2022} del INDEC (28 radios del departamento, 1.076,4 km²) y la \textbf{«Catastro Urbano y Rural de la Provincia de Salta»} de la Dirección General de Inmuebles, descargada del Geoportal de IDESA en septiembre de 2026 (6.219 parcelas del departamento). El organismo advierte que esta última es orientativa. Superficies calculadas en POSGAR 94, faja 3.


Se usaron además, tal como los publica el INDEC, los cuadros definitivos del Censo 2022 para la provincia de Salta (noviembre de 2023): el 1.17 (población por departamento, 2010 y 2022), el 2.17 (superficie y densidad) y los de condiciones habitacionales de la población 2.17.11, 3.17.11 y 6.17 (agua, desagüe y tenencia de la vivienda).


Las tasas de actividad, empleo y desocupación del capítulo \ref{cap:trabajo}, y las relaciones entre población y vivienda del capítulo \ref{cap:poblacion}, se obtuvieron \textbf{procesando el Censo Nacional 2010 ---cuestionario ampliado--- en Redatam}, la herramienta pública de CEPAL/CELADE con que se difunde la base censal. La consulta se hizo en \textbf{septiembre de 2026} y es reproducible: cualquiera puede repetirla y obtener los mismos valores. \textbf{El censo 2022 se procesó del mismo modo} y sus tasas figuran junto a las de 2010 en ese capítulo.

\textbf{Es, con las series hidrológicas y las mediciones sobre imagen, una de las fuentes de este libro que no son un documento sino un cálculo}, y por eso se declara la base, la fecha y el nivel de agregación. Cuando una misma consulta devuelve valores idénticos sobre bases distintas ---como ocurrió con las dos bases de 2022--- el capítulo lo consigna y no infiere nada de esa coincidencia.

\section{Documentos de vecinos difundidos por prensa}

El capítulo \ref{cap:resistencias} usa un manifiesto de Vecinos Autoconvocados difundido por \textit{Cuarto Poder} (16 de octubre de 2018) y por \textit{La Gaceta Salta} (18 de octubre de 2018). \textbf{Este libro no tiene el original}: lo que tiene son las dos notas, cotejadas letra por letra en septiembre de 2026.

Se usa porque \textbf{tres de sus afirmaciones verificables coinciden} con lo que el libro documenta por expediente, y porque acredita algo que ninguna otra fuente acredita: que hubo oposición organizada y cuándo.

\textbf{Su parte imputativa recibe otro tratamiento.} Una denuncia recogida por prensa no es un hecho probado: es una denuncia. El capítulo la consigna como tal, declara que no la verificó, y remite al protocolo de derecho a réplica antes de cualquier publicación. \textbf{Es la imputación personal más directa que aparece en este libro, y recibe el tratamiento más estricto.}

\section{Declaraciones de funcionarios en prensa}

El capítulo \ref{cap:vaqueros} usa declaraciones del director provincial de Vialidad Nacional publicadas por un medio especializado, \textit{Construar}, el 21 de agosto de 2026; el capítulo \ref{cap:prospectiva} las retoma. Lo que la nota afirma en su propia voz, sin comillas, se atribuye al medio y no al funcionario. \textbf{Una declaración no es un acto administrativo}: acredita lo que un funcionario dijo públicamente, no lo que el organismo resolvió, y no individualiza nada.

Este libro las usa sólo cuando el funcionario describe el estado de algo que le compete, las cita textualmente y pide, en el apéndice \ref{cap:pedidos}, el documento que las respaldaría. \textbf{Donde una declaración sea la única fuente de una afirmación, el capítulo lo dice.}

\section{Fuentes académicas}

\textbf{Este libro no hizo una búsqueda sistemática de bibliografía académica sobre la cuenca}: no revisó revistas, tesis ni repositorios en busca de trabajos sobre La Caldera, y usa sólo las cinco fuentes académicas que el último apartado de este apéndice enumera. Es una carencia y se declara como tal.

Entre ellas, la única que funda directamente una norma es el \textbf{resumen del informe técnico de revisión del Ordenamiento Territorial de Bosques Nativos}, publicado en el repositorio institucional del organismo nacional de ciencia, con sus autores identificados. Se usa en el capítulo \ref{cap:opacidad} \textbf{porque es el documento que fundó una ley que el libro ya analizaba}, no porque el libro haya hecho revisión bibliográfica.

\section{Un límite que se saldó}

Durante buena parte de este relevamiento, el Informe Final del \textbf{PDES 2030} fue la única fuente del libro leída sólo en parte: se habían consultado sus primeras ciento cincuenta páginas y no los capítulos de planeamiento urbano territorial, medio ambiente y turismo.

\textbf{Esa limitación se levantó}: el documento se recorrió completo, y el resultado ---que «La Caldera» aparece dos veces en ciento setenta páginas--- está en el capítulo \ref{cap:trabajo}.

\textbf{Se deja constancia de que el límite existió} porque durante ese lapso el libro afirmó, sobre una lectura parcial, algo que la lectura completa terminó confirmando. \textbf{Podría no haberlo hecho}, y la diferencia entre una cosa y otra no la decide el autor sino la lectura.

\section{Visores y portales cartográficos}

Este libro consultó el visor público de la Infraestructura de Datos Espaciales de la provincia
---IDESA--- y lo usa en el capítulo \ref{cap:opacidad}. \textbf{No es una fuente documental}: es
la observación del comportamiento de un sitio web en una fecha, y como tal se consigna, con
fecha y con la aclaración de qué se probó y qué no.

Lo mismo vale para cualquier consulta a un visor: acredita lo que el visor mostraba ese día, no
lo que el organismo tiene.

\section{Fuente genealógica}

Las fechas de nacimiento, matrimonio y muerte de Silvano Ignacio Murúa y de su padre (apéndice
\ref{cap:personas}) provienen de \textbf{Genealogía Familiar} (genealogiafamiliar.net), compilación
en línea que cita a su vez la desaparecida \textit{Genealogía Salteña} y registros de FamilySearch,
consultada en septiembre de 2026. \textbf{Es una fuente secundaria}: se usa para fechar, no para
probar dominio, y cada fecha debería cotejarse contra la partida del Registro Civil o el libro
parroquial. Su utilidad aquí es de control: las edades que da son compatibles con todas las
apariciones de los Murúa en el Boletín entre 1908 y 1944.

\section{Fuentes producidas en el territorio}

\textit{Nuestro tiempo redondo}, libro de la Comunidad Kolla Kondorwaira editado por la Municipalidad de La Caldera (2019), es la única fuente escrita por habitantes del departamento que este libro usa (capítulo \ref{cap:resistencias}). El \textbf{Plan Integral de Manejo} de la microcuenca (YSATI Ingeniería, 2024), con su diagnóstico social basado en encuestas anónimas, es la única que recoge testimonios de vecinos con método (capítulo \ref{cap:amparo}).

\section{Prensa}

Los medios con cobertura sostenida se enumeran en el capítulo \ref{cap:metodo}. Cada nota se cita donde se usa, con medio y fecha.

\section{Referencias teóricas}

El capítulo \ref{cap:prospectiva} discute el concepto de dispositivo en Michel Foucault (1977) y en Giorgio Agamben (2006). Son las únicas referencias teóricas del libro y se usan para fijar una definición, no como fuente de hechos.

\section{Lo que no está}

\begin{cautela}
Cuatro ausencias que el lector debe tener presentes, y que no son omisiones de este apéndice
sino del trabajo que lo precede.

\textbf{No hay entrevistas.} Ninguna. Todo lo que este libro dice sobre criterios y prácticas
lo infiere de documentos, y el capítulo \ref{cap:metodo} explica por qué eso alcanza para
algunas afirmaciones y no para otras.

\textbf{No hay revisión bibliográfica.} La literatura académica se usa de manera puntual: los trabajos sobre Vaqueros como caso de transformación periurbana y sobre la urbanización del Gran Salta (capítulos \ref{cap:metodo}, \ref{cap:fincas} y \ref{cap:tierra}), el trabajo de Polliotto, Lema, Alonso Maurizzio y Reyes (UCASAL, 2019) (capítulo \ref{cap:loteo}), el análisis del amparo del observatorio de la Universidad Nacional del Litoral (capítulo \ref{cap:amparo}) y el resumen del informe técnico del OTBN (capítulo \ref{cap:opacidad}). No se hizo una revisión sistemática de revistas, tesis y repositorios sobre producción de suelo, regulación urbana o conflictos hídricos. Es una decisión de alcance, no un descuido, pero tiene un costo: el libro no sabe qué de lo que describe es propio de este departamento y qué
es un patrón general.

\textbf{No hay archivo municipal.} Ni de La Caldera ni de Vaqueros. Las ordenanzas que se citan llegaron por el Boletín provincial, por transcripción de terceros, por el sitio web del municipio de La Caldera ---el Código de Ordenamiento Territorial y los partes de prensa--- o por el \textbf{Digesto Jurídico Virtual de la Municipalidad de Vaqueros}, consultado en septiembre de 2026, del que salen, entre otros textos, el reglamento de la Ordenanza 437, el Código de Edificación de 2026 y la Ordenanza 469/15 (apéndice \ref{cap:pedidos}).

\textbf{No hay fuentes orales, y la memoria comunitaria está en un solo libro.} Salvo \textit{Nuestro tiempo redondo} y las encuestas anónimas del Plan de Manejo, el departamento aparece en estas
páginas tal como el Estado lo escribió.
\end{cautela}
